% ==============================================
% About: Ficha resumen
% Description: Resumen técnico trabajo de grado
% ==============================================

%% MAKE FICHA SE PUEDE LLEVAR A OTRO LUGAR NO EN Thesis.cls
\newcommand{\makeficha}{
    \ \newpage
    \thispagestyle{empty}
    \begin{center}
    % \includegraphics[width=8cm]{./img/pujlogo.png}
    % \includegraphics[width=9.4cm]{./img/pujlogo2022.png}
    % \includegraphics[width=\textwidth]{./img/IngEl_FIC_PUJC.png}
    \includegraphics[width=15cm]{images/MSc/IngEl_FIC_PUJC.png}
    \end{center}
    %-----  CARTA PARA EL DIRECTOR DE CARRERA

    % \thispagestyle{empty}
    \begin{center}
     \bf{FICHA RESUMEN}
    \vspace{0.5cm}
    \bf{TRABAJO DE GRADO DE MAESTRÍA} \vspace{0.5cm}
    \end{center}

    %\noindent Agradecimientos. \\
    \vspace{0.5cm}

    % \noindent Dear\\
    \noindent {\bf TITULO:}  \titulovar \vspace{0.5cm}
    % \vspace{0.05cm}\\
    % \noindent {\bf Cédula:}   1.144'095.556\\
    % \vspace{0.1cm}\\
    \begin{enumerate}
        \item {\bf ÉNFASIS:} Ingeniería  -  Énfasis en electrónica.
        \item {\bf TIPO DE PROYECTO:} Tésis II -Proyecto de investigación.
        \item {\bf ÁREA DE TRABAJO:} Agropecuaria
        \item {\bf ESTUDIANTE:} \authorA
        \item {\bf CORREO ELECTRÓNICO:} \href{mailto:lfguevarag97@gmail.com}{lfguevarag97@gmail.com}
        \item {\bf DIRECCIÓN Y TELÉFONO:} Carrera 87 \#10 - 55, Multicentro 10, Apartamento 408. (+57) 305 -  379 80 30
        \item {\bf DIRECTOR:} \director
        \item {\bf VINCULACIÓN DEL DIRECTOR:} Profesor Titular, Director de carrera de Ingeniería Electrónica
        \item {\bf CORREO ELECTRÓNICO DEL DIRECTOR:}  \href{mailto:letobon@javerianacali.edu.co}{letobon@javerianacali.edu.co}
        \item {\bf CO-DIRECTOR(ES):} N/A
        \item {\bf GRUPO QUE LO AVALA:}  Servicio Nacional de Aprendizaje SENA, Seccional Popayán
        \item {\bf OTROS GRUPOS:} N/A
        \item {\bf PALABRAS CLAVE:} Modelación Dinámica, Ajuste de curvas, Producción ganadera, Regresión no lineal. Splines, Ecuaciones Diferenciales Ordinarias (EDO), Levenberg-Marquardt (LMA), Efectos no lineales mixtos (NLMEA).
        \pagebreak
        \item {\bf ODS QUE APLICA EL PROYECTO (AGENDA 2030):}
        \begin{itemize}
            \item ODS \#9 - Industria, innovación e infraestructura: Apoyar el desarrollo de tecnologías, la investigación y la innovación nacionales en los países en desarrollo, incluso garantizando un entorno normativo propicio a la diversificación industrial y la adición de valor a los productos básicos, entre otras cosas (9.b).
            \item ODS \#12 - Producción y consumo responsables: Ayudar a los países en desarrollo a fortalecer su capacidad científica y tecnológica para avanzar hacia modalidades de consumo y producción más sostenibles (12.a).
        \end{itemize}

        % \item {\bf CODIGOS UNESCO - CIENCIA Y TECNOLOGÍA:} 1203.26, 1206.07, 1206.12, 3104.90
        \item  \textbf{FECHA DE INICIO:} 1 de Febrero de 2023.% - {\bf DURACIÓN ESTIMADA:} 10 meses
        \item  \textbf{RESUMEN:} Ver siguiente página.

    \end{enumerate}
    \cleardoublepage
}
 % End of MakeFicha Tecnica
